% ============================================================
%  Chapter 4 -- Continuity and Compactness, 4.13-4.21
% ============================================================
\rsection{Continuity and Compactness}

\begin{signpost}[Chapter 2's investment matures]
Five theorems, one engine. 4.14 shows continuous images of compact sets are
compact --- a one-paragraph proof, because 4.8 lets the compactness of the
domain be spent directly on a pulled-back cover. Everything else in the
section is 4.14 plus a Chapter 2 fact: boundedness (4.15, via Heine--Borel),
attainment of extrema (4.16, via 2.28), continuity of inverse maps (4.17),
and uniform continuity (4.19, where compactness converts infinitely many
local $\delta$'s into one global one --- the finite-minimum move yet again,
and Rudin says so in the proof).

\smallskip
Then the audit: 4.20--4.21 dismantle each theorem when compactness is
removed, one counterexample per conclusion. The section is a matched pair ---
what compactness buys, and the invoice showing nothing was free.
\end{signpost}

\ritem{4.13}{Definition}
A mapping $\mathbf{f}$ of a set $E$ into $\R^k$ is said to be \emph{bounded}
if there is a real number $M$ such that $|\mathbf{f}(x)| \le M$ for all
$x \in E$.

\ritem{4.14}{Theorem}
\emph{Suppose $f$ is a continuous mapping of a compact metric space $X$ into
a metric space $Y$. Then $f(X)$ is compact.}

\begin{proof}
Let $\{V_\alpha\}$ be an open cover of $f(X)$. Since $f$ is continuous,
Theorem 4.8 shows that each of the sets $f^{-1}(V_\alpha)$ is open. Since $X$
is compact, there are finitely many indices, say $\alpha_1,\dots,\alpha_n$,
such that
\begin{equation}
  X \subset f^{-1}(V_{\alpha_1}) \cup \cdots \cup f^{-1}(V_{\alpha_n}).
  \tag{4.12}
\end{equation}
Since $f(f^{-1}(E)) \subset E$ for every $E \subset Y$, (4.12) implies that
\begin{equation}
  f(X) \subset V_{\alpha_1} \cup \cdots \cup V_{\alpha_n}. \tag{4.13}
\end{equation}
This completes the proof.
\end{proof}

\emph{Note:} We have used the relation $f(f^{-1}(E)) \subset E$, valid for
$E \subset Y$. If $E \subset X$, then $f^{-1}(f(E)) \supset E$; equality need
not hold in either case.

\begin{center}
\begin{tikzpicture}[x=1mm, y=1mm]
  \node[rmbox, text width=28mm] (f0) at (0,20)
    {open cover $\{V_\alpha\}$\\of $f(X)$, in $Y$};
  \node[rmbox, text width=29mm] (f1) at (43,20)
    {pull back: $f^{-1}(V_\alpha)$\\open cover of $X$};
  \node[rmbox, text width=23mm] (f2) at (89,20)
    {finitely many\\cover $X$};
  \draw[rmarrow] (f0) -- node[lbl, anchor=south, inner sep=2.5pt] {4.8} (f1);
  \draw[rmarrow] (f1) -- node[lbl, anchor=south, inner sep=2.5pt,
    align=center] {$X$\\compact} (f2);
  \node[rmbox, text width=33mm] (f3) at (43,0)
    {push forward: the same\\finitely many $V_\alpha$ cover $f(X)$};
  \draw[rmarrow] (f2) -- node[lbl, anchor=north west, inner sep=2.5pt]
    {$f(f^{-1}(E)) \subset E$} (f3);
\end{tikzpicture}
\end{center}
\figcap{Theorem 4.14 as a flow diagram: pull the cover back, spend
compactness, push the finite subcover forward. Each arrow is one theorem, and
the round trip works because $f$ is defined on \emph{all} of $X$ --- the
push-forward inclusion needs nothing else.}

We shall now deduce some consequences of Theorem 4.14.

\ritem{4.15}{Theorem}
\emph{If $\mathbf{f}$ is a continuous mapping of a compact metric space $X$
into $\R^k$, then $\mathbf{f}(X)$ is closed and bounded. Thus, $\mathbf{f}$
is bounded.}

This follows from Theorem 2.41. The result is particularly important when $f$
is real:

\ritem{4.16}{Theorem}
\emph{Suppose $f$ is a continuous real function on a compact metric space
$X$, and}
\begin{equation}
  M = \sup_{p\in X} f(p), \qquad m = \inf_{p\in X} f(p). \tag{4.14}
\end{equation}
\emph{Then there exist points $p, q \in X$ such that $f(p) = M$ and
$f(q) = m$.}

The notation in (4.14) means that $M$ is the least upper bound of the set of
all numbers $f(p)$, where $p$ ranges over $X$, and that $m$ is the greatest
lower bound of this set of numbers.

The conclusion may also be stated as follows: \emph{There exist points $p$
and $q$ in $X$ such that $f(q) \le f(x) \le f(p)$ for all $x \in X$}; that
is, $f$ attains its maximum (at $p$) and its minimum (at $q$).

\begin{proof}
By Theorem 4.15, $f(X)$ is a closed and bounded set of real numbers; hence
$f(X)$ contains
\[ M = \sup f(X) \qquad\text{and}\qquad m = \inf f(X), \]
by Theorem 2.28.\kn{The extreme value theorem of calculus, in three
citations: 4.14 makes the image compact, 2.41 reads that as closed and
bounded, 2.28 puts the supremum of a closed bounded set \emph{inside} it.
Every hard step was prepaid in Chapter 2 --- exactly as the note at 2.28
promised. Notice the theorem gives no way to \emph{find} $p$ and $q$; pure
existence, from compactness alone.}
\end{proof}

\ritem{4.17}{Theorem}
\emph{Suppose $f$ is a continuous 1-1 mapping of a compact metric space $X$
onto a metric space $Y$. Then the inverse mapping $f^{-1}$ defined on $Y$ by}
\[ f^{-1}(f(x)) = x \qquad (x \in X) \]
\emph{is a continuous mapping of $Y$ onto $X$.}

\begin{proof}
Applying Theorem 4.8 to $f^{-1}$ in place of $f$, we see that it suffices to
prove that $f(V)$ is an open set in $Y$ for every open set $V$ in $X$. Fix
such a set $V$.

The complement $V^c$ of $V$ is closed in $X$, hence compact (Theorem 2.35);
hence $f(V^c)$ is a compact subset of $Y$ (Theorem 4.14) and so is closed in
$Y$ (Theorem 2.34). Since $f$ is one-to-one and onto, $f(V)$ is the
complement of $f(V^c)$. Hence $f(V)$ is open.\kn{A ricochet proof: to open
$f(V)$, close its complement. The chain closed $\to$ compact $\to$ compact
image $\to$ closed converts domain-compactness into the openness 4.8 demands
of $f^{-1}$ --- and each link is a Chapter 2 theorem (2.35, 4.14, 2.34).
Bijectivity enters exactly once, to say image-of-complement is
complement-of-image. Example 4.21 shows the whole chain snaps without
compactness.}
\end{proof}

\ritem{4.18}{Definition}
Let $f$ be a mapping of a metric space $X$ into a metric space $Y$. We say
that $f$ is \emph{uniformly continuous} on $X$ if for every $\varepsilon > 0$
there exists $\delta > 0$ such that
\begin{equation}
  d_Y(f(p),f(q)) < \varepsilon \tag{4.15}
\end{equation}
for all $p$ and $q$ in $X$ for which $d_X(p,q) < \delta$.

Let us consider the differences between the concepts of continuity and of
uniform continuity. First, uniform continuity is a property of a function on
a set, whereas continuity can be defined at a single point. To ask whether a
given function is uniformly continuous at a certain point is meaningless.
Second, if $f$ is continuous on $X$, then it is possible to find, for each
$\varepsilon > 0$ and \emph{for each point $p$ of $X$}, a number $\delta > 0$
having the property specified in Definition 4.5. This $\delta$ depends on
$\varepsilon$ \emph{and on $p$}. If $f$ is, however, uniformly continuous on
$X$, then it is possible, for each $\varepsilon > 0$, to find \emph{one}
number $\delta > 0$ which will do for \emph{all} points $p$ of
$X$.\kw{Quantifier order is the entire difference: continuity is
$\forall p\,\forall\varepsilon\,\exists\delta$, uniform continuity is
$\forall\varepsilon\,\exists\delta\,\forall p$. Moving $\forall p$ past
$\exists\delta$ strips $\delta$ of its dependence on $p$. This is the
book's sharpest lesson in reading quantifiers, and the figure below is the
picture to attach to it.}

Evidently, every uniformly continuous function is continuous. That the two
concepts are equivalent on compact sets follows from the next theorem.

\begin{center}
\begin{tikzpicture}[x=1mm, y=1mm]
  % curve y = 1/x-like steepening on (0, ...]
  \draw[axis] (-2,0) -- (94,0);
  \draw[axis] (0,-2) -- (0,34);
  \draw[thin] plot[smooth, tension=0.8]
    coordinates {(4,32) (8,17) (14,9.5) (24,5.5) (44,3.0) (90,1.6)};
  \node[lbl, anchor=west] at (68,6.5) {$f(x)=1/x$ on $(0,1]$};
  % same epsilon window at two points, different deltas
  \draw[guide] (8,-4) -- (8,17);
  \draw[thin, {Stealth[length=3pt]}-{Stealth[length=3pt]}] (6.9,-4) -- (9.1,-4);
  \node[lbl, anchor=north] at (8,-5.5) {tiny $\delta$ needed};
  \draw[guide] (44,-4) -- (44,3.0);
  \draw[thin, {Stealth[length=3pt]}-{Stealth[length=3pt]}] (36,-4) -- (52,-4);
  \node[lbl, anchor=north] at (44,-5.5) {generous $\delta$ suffices};
\end{tikzpicture}
\end{center}
\figcap{Continuity that is not uniform: the same $\varepsilon$ of vertical
tolerance costs ever-smaller $\delta$ as the curve steepens toward the
missing endpoint. No single $\delta$ serves every point --- Theorem 4.20(c)
makes this the general phenomenon on bounded non-compact sets.}

\ritem{4.19}{Theorem}
\emph{Let $f$ be a continuous mapping of a compact metric space $X$ into a
metric space $Y$. Then $f$ is uniformly continuous on $X$.}

\begin{proof}
Let $\varepsilon > 0$ be given. Since $f$ is continuous, we can associate to
each point $p \in X$ a positive number $\phi(p)$ such that
\begin{equation}
  q \in X,\ d_X(p,q) < \phi(p) \quad\text{implies}\quad
  d_Y(f(p),f(q)) < \frac{\varepsilon}{2}. \tag{4.16}
\end{equation}
Let $J(p)$ be the set of all $q \in X$ for which
\begin{equation}
  d_X(p,q) < \tfrac12\phi(p). \tag{4.17}
\end{equation}
Since $p \in J(p)$, the collection of all sets $J(p)$ is an open cover of
$X$; since $X$ is compact, there is a finite set of points $p_1,\dots,p_n$ in
$X$, such that
\begin{equation}
  X \subset J(p_1) \cup \cdots \cup J(p_n). \tag{4.18}
\end{equation}
We put
\begin{equation}
  \delta = \tfrac12\min\{\phi(p_1),\dots,\phi(p_n)\}. \tag{4.19}
\end{equation}
Then $\delta > 0$. (This is one point where the finiteness of the covering,
inherent in the definition of compactness, is essential. The minimum of a
finite set of positive numbers is positive, whereas the inf of an infinite
set of positive numbers may very well be $0$.)

Now let $q$ and $p$ be points of $X$ such that $d_X(p,q) < \delta$. By
(4.18), there is an integer $m$, $1 \le m \le n$, such that $p \in J(p_m)$;
hence
\begin{equation}
  d_X(p,p_m) < \tfrac12\phi(p_m), \tag{4.20}
\end{equation}
and we also have
\[ d_X(q,p_m) \le d_X(p,q) + d_X(p,p_m)
   < \delta + \tfrac12\phi(p_m) \le \phi(p_m). \]
Finally, (4.16) shows that therefore
\[ d_Y(f(p),f(q)) \le d_Y(f(p),f(p_m)) + d_Y(f(q),f(p_m)) < \varepsilon. \]
This completes the proof.
\end{proof}

An alternative proof is sketched in Exercise 4.10.

\begin{proofwalk}[How the proof flows]
  \item \textbf{Name the enemy.} Continuity hands out a private
        $\delta$-radius $\phi(p)$ at each point; uniformity demands one
        public $\delta$. The natural attempt --- take
        $\inf_p \phi(p)$ --- fails because an infimum of infinitely many
        positive numbers can be $0$ (Rudin says exactly this in the
        parenthesis at (4.19)). The proof exists to make that infimum
        finite.
  \item \textbf{Turn the private radii into a cover.} Each point owns the
        ball $J(p)$ of \emph{half} its radius; together they cover $X$.
        Compactness reduces to finitely many centers $p_1,\dots,p_n$ --- and
        now the minimum in (4.19) is over a finite set, hence positive. This
        is the local-to-global conversion promised in Chapter 2's signpost,
        performed literally.
  \item \textbf{Why half?} The safety margin is what makes the public
        $\delta$ work for pairs. A point $p$ lands in some $J(p_m)$, within
        $\tfrac12\phi(p_m)$ of the center; its partner $q$ is within
        $\delta \le \tfrac12\phi(p_m)$ more, so \emph{both} lie inside the
        full radius $\phi(p_m)$. Without the halving, $q$ could fall just
        outside.
  \item \textbf{Close through the center.} Neither $p$ nor $q$ is compared
        with the other directly --- both are compared with $p_m$, each
        within $\varepsilon/2$, and the triangle inequality adds the halves.
        That is why (4.16) was set up with $\varepsilon/2$: working
        backwards, pattern 4.
  \item \textbf{Audit.} Compactness spent once (at the finite subcover),
        continuity spent once per point (to define $\phi$), triangle
        inequality twice. The halving trick --- shrink each local radius so
        that overlapping neighbours still see each other --- recurs in
        Chapter 7 and in every partition-of-unity argument (Chapter 10).
\end{proofwalk}

We now proceed to show that compactness is essential in the hypotheses of
Theorems 4.14, 4.15, 4.16, and 4.19.

\ritem{4.20}{Theorem}
\emph{Let $E$ be a noncompact set in $\R^1$. Then}
\begin{enumerate}[label=(\alph*), leftmargin=2.2em, itemsep=1pt, topsep=3pt]
  \item \emph{there exists a continuous function on $E$ which is not
        bounded;}
  \item \emph{there exists a continuous and bounded function on $E$ which
        has no maximum.}
\end{enumerate}
\emph{If, in addition, $E$ is bounded, then}
\begin{enumerate}[label=(\alph*), leftmargin=2.2em, itemsep=1pt, topsep=3pt,
                  start=3]
  \item \emph{there exists a continuous function on $E$ which is not
        uniformly continuous.}
\end{enumerate}

\begin{proof}
Suppose first that $E$ is bounded, so that there exists a limit point $x_0$
of $E$ which is not a point of $E$.\kn{Here 2.41 is being run in reverse:
noncompact-but-bounded forces not-closed, which hands the proof its weapon
--- a missing limit point $x_0$. Every counterexample in the bounded case is
built by blowing up at $x_0$; the theorem is really about how continuity on
$E$ cannot see the hole just outside $E$.} Consider
\begin{equation}
  f(x) = \frac{1}{x-x_0} \qquad (x \in E). \tag{4.21}
\end{equation}
This is continuous on $E$ (Theorem 4.9), but evidently unbounded. To see that
(4.21) is not uniformly continuous, let $\varepsilon > 0$ and $\delta > 0$ be
arbitrary, and choose a point $x \in E$ such that $|x-x_0| < \delta$. Taking
$t$ close enough to $x_0$, we can then make the difference $|f(t)-f(x)|$
greater than $\varepsilon$, although $|t-x| < \delta$. Since this is true for
every $\delta > 0$, $f$ is not uniformly continuous on $E$.

The function $g$ given by
\begin{equation}
  g(x) = \frac{1}{1+(x-x_0)^2} \qquad (x \in E) \tag{4.22}
\end{equation}
is continuous on $E$ and is bounded, since $0 < g(x) < 1$. It is clear that
\[ \sup_{x\in E} g(x) = 1, \]
whereas $g(x) < 1$ for all $x \in E$. Thus $g$ has no maximum on $E$.

Having proved the theorem for bounded sets $E$, let us now suppose that $E$
is unbounded. Then $f(x) = x$ establishes (a), whereas
\begin{equation}
  h(x) = \frac{x^2}{1+x^2} \qquad (x \in E) \tag{4.23}
\end{equation}
establishes (b), since
\[ \sup_{x\in E} h(x) = 1 \]
and $h(x) < 1$ for all $x \in E$.

Assertion (c) would be false if boundedness were omitted from the
hypotheses. For, let $E$ be the set of all integers. Then every function
defined on $E$ is uniformly continuous on $E$. To see this, we need merely
take $\delta < 1$ in Definition 4.18.\kd{So the converse pattern is not
clean: noncompactness kills boundedness and attainment outright, but uniform
continuity survives on $\Z$ --- an unbounded, closed, noncompact set where
every function is uniformly continuous because points cannot get close. The
moral: (a) and (b) fail through \emph{escape} (to a missing point or to
infinity), (c) fails only through \emph{crowding} against a missing point.
Different failure modes, different hypotheses. Exercise 4.8 probes the same
boundary.}
\end{proof}

We conclude this section by showing that compactness is also essential in
Theorem 4.17.

\ritem{4.21}{Example}
Let $X$ be the half-open interval $[0,2\pi)$ on the real line, and let
$\mathbf{f}$ be the mapping of $X$ onto the circle $Y$ consisting of all
points whose distance from the origin is $1$, given by
\begin{equation}
  \mathbf{f}(t) = (\cos t, \sin t) \qquad (0 \le t < 2\pi). \tag{4.24}
\end{equation}
The continuity of the trigonometric functions cosine and sine, as well as
their periodicity properties, will be established in Chap.\ 8. These results
show that $\mathbf{f}$ is a continuous 1-1 mapping of $X$ onto $Y$.

However, the inverse mapping (which exists, since $\mathbf{f}$ is one-to-one
and onto) fails to be continuous at the point $(1,0) = \mathbf{f}(0)$. Of
course, $X$ is not compact in this example. (It may be of interest to observe
that $\mathbf{f}^{-1}$ fails to be continuous in spite of the fact that $Y$
\emph{is} compact!)

\begin{center}
\begin{tikzpicture}[x=1mm, y=1mm]
  % the half-open interval
  \draw[axis] (-2,0) -- (46,0);
  \draw[seg] (2,0) -- (40,0);
  \node[ptfill, minimum size=2.6pt] at (2,0) {};
  \node[ptopen] at (40,0) {};
  \node[mlbl, anchor=north] at (2,-2) {$0$};
  \node[mlbl, anchor=north] at (40,-2) {$2\pi$};
  \node[lbl, anchor=north] at (21,-6) {$X=[0,2\pi)$: the seam is open};
  % wrap arrow
  \draw[vec] (48,4) .. controls (54,8) .. (58,4);
  \node[mlbl, anchor=south] at (53,8) {$\mathbf{f}$};
  % the circle
  \draw[thin] (74,4) circle (11mm);
  \node[ptfill, minimum size=2.6pt] at (85,4) {};
  \node[mlbl, anchor=west] at (86.5,4) {$(1,0)$};
  \draw[thin, -{Stealth[length=3pt]}] (84.2,7.8) .. controls (80,11) .. (75,10.6);
  \draw[thin, -{Stealth[length=3pt]}] (84.2,0.2) .. controls (80,-3) .. (75,-2.6);
  \node[lbl, anchor=north] at (74,-9) {$Y$: both ends of $X$ arrive at $(1,0)$};
\end{tikzpicture}
\end{center}
\figcap{Example 4.21: wrapping the half-open interval onto the circle glues
the missing endpoint to $0$. Points just below $2\pi$ land next to $(1,0)$,
but $\mathbf{f}^{-1}$ throws them back near $2\pi$ --- a jump of almost
$2\pi$ across the seam. Continuity of the inverse is a property the
\emph{domain's} compactness must fund; the target's cannot.}
